
% !TeX root = ../main.tex
\chapter{Discussion}\label{chapter:discussion}

This chapter interprets the relationships between challenges and solutions and qualifies the conclusions supported by the reviewed literature.

\section{Principal Findings}

The analysis identified ten challenges across five themes and 34 solution mechanisms. Table~\ref{table:challenges-and-solutions-overview} maps the relationships between these mechanisms and challenges as reported by the source publications. The mapping does not establish the applicability of either a challenge or a solution mechanism to a target configuration.

The challenge themes distinguish several consequences associated in the literature with distributing different search functions. The reviewed publications connect client-side search workloads to resource constraints, distributed candidate evaluation to dependence on local data, and decentralized search-state updates to coordination requirements. These synthesized relationships describe the reviewed evidence; they do not predict whether a challenge will occur in a system not examined by those publications.

The publications also report interactions between solution mechanisms. Assigning capacity-compatible subnets reduces client workloads, but causes operations available only to capable clients to receive fewer updates. Exposure-aware aggregation is then used to prevent unequal training from distorting the shared state~\cite{fedoras_2022,fgfl_2024,supernet_trade-off_fednas_2025}. Similarly, specialized architectures accommodate different deployment requirements while removing the parameter alignment required for ordinary aggregation. The reviewed methods use structural aggregation or distillation to restore collaboration under this condition~\cite{multi-granular_fednas_2023,f2nas_2024,hapfnas_2025}. The map records these reported dependencies but does not demonstrate that combining the linked mechanisms is sufficient in another configuration.

The search space creates a recurring trade-off between architectural coverage and coordination. Restricting it to inexpensive, structurally aligned components lowers client workloads and enables direct aggregation, but may exclude architectures that better fit particular data or devices. Elastic or personalized spaces cover more deployment conditions, but increase search effort and can produce unequal parameter exposure or incompatible local states. The search space thus serves both as the set of possible architectures and as a coordination contract between clients.

Brief training is the predominant performance-estimation strategy in the reviewed literature. Within that broad category, the publications differ in the workload assigned to clients, the relationship between their search states, and whether the search produces a shared or specialized architecture. Broad search-strategy categories alone can conceal these differences, so Chapter~\ref{chapter:nas_in_fl_ch_and_sol} retains more specific NAS design choices and FL system characteristics as evidence context rather than as a basis for inferring applicability.

\section{Implications for Practice and Research}

The synthesis provides an index of solution mechanisms that publications present in response to particular challenges. Practitioners can use the map to locate relevant mechanisms and their source evaluations, but the map is not a procedure for selecting a mechanism for a target system. Such a choice requires direct evidence for that system's operating conditions and cannot be inferred from a matching description alone.

For research, the synthesis complements method-level taxonomies by treating challenges and solution mechanisms as separate units. This exposes relationships that remain hidden when several mechanisms are presented as one NAS-in-FL method. It also provides a common vocabulary for comparing reported mechanisms and their evidence contexts. Future publications can support such comparisons by reporting the FL system characteristics, client workloads, exchanged search state, search space, and performance-estimation strategy used in their evaluations.

\section{Limitations}

The review depends on what the selected publications report. It can include only challenges, solution mechanisms, relationships, prerequisites, and trade-offs described with enough detail to meet the inclusion criteria. Operational problems that have not been published or that receive only general descriptions may therefore be absent. The Scopus search and citation searching reduce this risk but do not establish that the corpus covers every relevant publication.

The thematic analysis also requires judgment when concepts are named, separated, merged, and mapped~\cite{using_thematic_analysis_2006}. Explicit definitions and an audit trail make these decisions transparent, but do not eliminate interpretive subjectivity. Moreover, few publications isolate the effect of one FL characteristic or solution mechanism. The resulting mappings organize reported relationships; they are not independent estimates of causal effects, applicability, or comparative effectiveness.

Evaluations differ in their datasets, search spaces, client populations, participation, resource models, and training budgets. These differences prevent a quantitative ranking of the mechanisms. Incomplete descriptions of FL systems further prevent comparison across settings because an unreported condition cannot be distinguished from one that was absent. Publication frequency therefore indicates attention within the corpus, not practical prevalence, severity, or solution quality.

\section{Future Research}

Future evaluations should test individual mechanisms and selected combinations under shared configurations. Varying one FL condition or NAS property at a time would provide stronger evidence for the relationships identified in this review. Such evaluations should separate search computation, final-model training, communication, wall-clock time, candidate-ranking fidelity, and final architecture quality.

Scale and participation require particular attention. Experiments should vary the number of clients together with partial participation, dropout, latency, and capability-dependent selection, since a larger simulated population alone does not reproduce cross-device dynamics. Publications should also distinguish costs measured on client hardware from those simulated on centralized infrastructure.

A shared reporting scheme would make results easier to compare. At minimum, publications should describe client resources, network conditions, data distribution, participation, federation size, trust assumptions, client search workloads, exchanged search state, search space, and final-model development. Research should also examine sparsely studied areas of the challenge--solution map and test whether low-cost performance estimators from centralized NAS remain reliable when candidate evidence is distributed across heterogeneous client data~\cite{nas_in_fl_2026}.
